\section{Introduction}
\label{sec:intro}

Network intrusion detection, botnet identification, and IoT malware analysis
are canonical cybersecurity tasks that operate over inherently
\emph{graph-structured} data: traffic flows between hosts form edges, devices
form nodes, and attack campaigns manifest as anomalous subgraph patterns.
The relational structure is not incidental. Lateral movement is defined by which
hosts contact which others; command-and-control beaconing is defined by
periodicity in a communication pattern; a distributed denial-of-service campaign
is defined by convergence of many sources on one target. None of these is
visible in an individual flow record considered in isolation, which is why graph
representations have displaced purely tabular ones for these tasks.
Graph Neural Networks (GNNs) have accordingly emerged as powerful classifiers,
achieving high accuracy on benchmarks such as CIC-IDS-2017
\citep{cicids2017}, UNSW-NB15 \citep{unswNB15}, and NF-BoT-IoT
\citep{nfbotiot} — yet they suffer from two critical, simultaneously
unresolved weaknesses.

\paragraph{W1 — Epistemic uncertainty from noisy and ambiguous features.}
Network telemetry is inherently imprecise: packet-level features exhibit
quantisation noise, flow statistics are aggregated over variable windows,
and labelling of attack vs.\ benign traffic is often soft or partially
overlapping \citep{graphids}.
The imprecision is not measurement error to be averaged away but a property of
the domain. A long-lived encrypted session may be a backup job or an
exfiltration channel; the flow-level evidence is frequently compatible with
both, and the ground-truth labels supplied with public benchmarks are themselves
the product of heuristic post-hoc attribution.
Standard GNN classifiers produce crisp binary decisions that discard this
inherent fuzziness, leading to brittle predictions and inflated
false-alarm rates in production SOC deployments.

\paragraph{W2 — Absence of formal coverage guarantees.}
Neither standard GNNs nor fuzzy-neural extensions provide
\emph{distribution-free, finite-sample guarantees} on their prediction sets.
A security analyst receiving an alert has no principled bound on the
probability that the true attack class lies within the model's output —
making risk quantification impossible for compliance-sensitive
(e.g., GDPR, NIS2) environments \citep{cp_survey}.
Softmax confidence is not a substitute: modern networks are systematically
overconfident, and a reported probability of $0.99$ carries no guarantee about
the frequency with which such predictions are correct.

Nor do the standard remedies close the gap. Post-hoc calibration methods such as
temperature scaling adjust confidence to match average accuracy on a held-out
set, but supply no finite-sample guarantee for an individual prediction and are
themselves fitted under the assumption that deployment resembles validation.
Bayesian neural networks and deep ensembles produce a predictive distribution
whose spread is often taken as an uncertainty estimate, yet the resulting
credible intervals are valid only if the prior and likelihood are correctly
specified — an assumption no one is in a position to defend for adversarial
network traffic — and ensembles multiply training and inference cost by the
number of members. Conformal prediction is attractive precisely because it
assumes none of this: it wraps an arbitrary, possibly misspecified predictor and
returns a guarantee that holds in finite samples for any data distribution,
requiring only exchangeability between calibration and test data. In a domain
where the data-generating process is partly under an adversary's control, an
assumption-light guarantee is worth more than a sharper one resting on
conditions that cannot be checked.

Transferring that guarantee to graphs is not immediate. In transductive node
classification the prediction for one node depends, through message passing, on
the features and labels of its neighbours, so calibration and test scores are
statistically coupled and the i.i.d.\ reasoning that motivates split conformal
prediction no longer applies directly. Validity must be re-established from a
permutation-invariance argument rather than assumed \citep{cfgnn}, and any
modification to the scoring function — score smoothing across edges, for
instance — risks weakening the very symmetry the argument relies on.

These two weaknesses compound one another. A model that discards the ambiguity
in its input will express unwarranted confidence in its output, and a model that
offers no calibrated statement about that confidence gives the analyst no basis
on which to triage. What an operator needs from an alert is not a single label
but a defensible statement of the form ``the true class lies in this set with
probability at least $1-\alpha$'', together with a human-readable account of
which features drove the assessment.

We argue further that a \emph{marginal} guarantee is insufficient for this
setting. A detector can satisfy a $90\%$ coverage target averaged over all
traffic while systematically under-covering a rare attack family, because the
dominant benign class determines the average. Averaging is the wrong summary
for a security operations centre: an alert concerning a rare family is not made
trustworthy by the detector being well calibrated on benign traffic. The
guarantee must be conditional on the region of the graph a host inhabits.

\subsection{Research Gaps}

Significant progress has been made independently along two complementary
lines: (a) \emph{conformal prediction for GNNs}
\citep{cfgnn,daps,clarkson23,rrgnn} and (b) \emph{fuzzy GNN architectures}
\citep{flgnn,fgat,fgnn_survey}.
The first supplies distribution-free guarantees but operates on opaque softmax
scores and yields no account of \emph{why} a set was returned; the second
supplies inspectable rule bases but no statement about how often the true class
is captured, since a fuzzy membership value carries no frequentist
interpretation. The survey of \citet{fgnn_survey} identifies exactly this
absence of guarantees as an open problem spanning the whole fuzzy GNN
literature.
Their explicit fusion remains an open problem, and it is not merely a matter of
composition: a fuzzy architecture produces a graded membership vector rather
than a probability, so the nonconformity score that mediates between the model
and the conformal layer must itself be defined in fuzzy terms.
Table~\ref{tab:gaps} positions FC-GNN against the five closest prior methods.

\begin{table}[t]
\caption{Research gap: no prior work simultaneously provides coverage
         guarantees, fuzzy aggregation, and cybersecurity validation.}
\label{tab:gaps}
\small
\begin{tabular}{lccc}
\toprule
Method & Coverage & Fuzzy & Cyb.\ Bench. \\
\midrule
CF-GNN \citep{cfgnn}        & \checkmark & $\times$ & $\times$ \\
RR-GNN \citep{rrgnn}        & \checkmark & $\times$ & $\times$ \\
FL-GNN \citep{flgnn}        & $\times$ & \checkmark & $\times$ \\
FGAT   \citep{fgat}         & $\times$ & \checkmark & $\times$ \\
GraphIDS \citep{graphids}   & $\times$ & $\times$ & \checkmark \\
\midrule
\textbf{FC-GNN (ours)}      & \checkmark & \checkmark & \checkmark \\
\bottomrule
\end{tabular}
\end{table}

\subsection{Contributions}

This paper makes four contributions:

\begin{enumerate}
\item \textbf{FC-GNN architecture.}  A novel GNN in which each
  message-passing layer incorporates a fuzzy-rule aggregation module that
  outputs a fuzzy membership distribution over classes per node, rather than
  a crisp softmax logit (Section~\ref{sec:method}).
  Neighbourhood aggregation is performed by rule firing under a product T-norm
  rather than by a fixed permutation-invariant operator, so the quantity
  propagated between nodes is the degree to which each interpretable rule is
  satisfied.

\item \textbf{Fuzzy Mondrian Conformal Prediction (FMCP).}  A principled
  extension of Graph-Structured Mondrian CP \citep{rrgnn} to fuzzy
  nonconformity scores defined via a Sugeno-type fuzzy integral, producing
  calibrated prediction sets with per-cluster conditional coverage guarantees
  (Section~\ref{sec:fmcp}).
  The Sugeno integral is the natural aggregation operator here because it
  respects the ordinal structure of graded membership without presuming the
  additivity that a probabilistic treatment would impose.

\item \textbf{Coverage theorem.}  A finite-sample theorem proving that FC-GNN's
  prediction sets achieve at least $(1-\alpha)$ marginal coverage and
  conditional coverage within graph communities, with a slack term
  $\delta(|Q_m|) = O(|Q_m|^{-1/2})$ (Section~\ref{sec:theorem}).
  We state the exchangeability and community-stability conditions the result
  requires explicitly, and identify where each fails in deployment.

\item \textbf{Cybersecurity benchmarks.}  Systematic evaluation on seven
  publicly available datasets spanning intrusion detection, botnet detection,
  and IoT malware, with temporal chronological splits and full ablation
  studies (Section~\ref{sec:experiments}).
  Splits are chronological rather than random, so calibration and test data are
  separated in time; this is a deliberately harsher protocol than the random
  splits usual in the CP-GNN literature, and it is what exposes the difference
  between marginal and conditional calibration.
\end{enumerate}

\subsection{Summary of Findings}

Three results are worth stating in advance, including one that does not favour
the proposed method.
First, the distinction between marginal and conditional calibration is
empirically large, not a technical nicety: methods that smooth conformity scores
across the graph achieve competitive average set sizes while failing the nominal
coverage target on most datasets under chronological splits, whereas
community-conditional calibration holds it.
Second, evaluation protocol dominates apparent performance. Under the random
splits customary in the CP-GNN literature, calibration and test points are drawn
from the same temporal regime by construction, and differences between methods
largely disappear; the chronological protocol we adopt separates them in time as
a deployment would, and it is what makes the differences visible.
Third, FC-GNN does not dominate on point-prediction accuracy, and we do not
claim that it does. On the most severely imbalanced benchmark its Macro-F1 is
poor in absolute terms, and its advantage over the strongest Mondrian baseline
in conditional coverage is real but modest. What it provides that the
alternatives do not is a conditional guarantee together with a rule-level
account of the decision, and we regard the accuracy cost as the price of that
combination rather than as an incidental shortcoming.

\subsection{Paper Structure}

Section~\ref{sec:related} reviews related work.
Section~\ref{sec:method} covers necessary preliminaries
(Section~\ref{sec:prelim}) and presents the FC-GNN architecture and the
FMCP calibration procedure.
Section~\ref{sec:experiments} reports empirical results and analyses the
extracted fuzzy rules (Section~\ref{sec:interpretability}).
Section~\ref{sec:discussion} discusses limitations.
Section~\ref{sec:conclusion} concludes.
